\section{The research target and its antecedents}
\label{sec:scope}

\subsection{An explicit unresolved coefficient problem}

The inspected repository is pinned to commit
\begin{center}\small\pin.\end{center}
The user-specified Polylogarithms and Transseries directories and all
three archives present in \path{docs/incoming} at that commit were
inventoried. The relevant integrated chapters and predecessor reports
were read directly. This is a focused audit of the claims used here,
not a review of the entire repository.

Question Q1 in the boundary-completion report \cite{repo-boundary}
asks for a coefficient theorem with
\begin{equation}\label{eq:window}
 L=\log n,\qquad s=3+\eps,\qquad
 \lambda=\frac{\eps L}{2}
\end{equation}
bounded, while the feedback distance from criticality varies on its
natural local-limit scale. That report proves a convergent joint
\emph{inverse} chart but keeps $s$ fixed in its coefficient theorem.
Older linear-tail reports provide leading moving-order limits and a
first central correction under eventually exact power-tail hypotheses
\cite{repo-gaussian,repo-confluent}. The quadratic outer function changes
infinitely many Taylor coefficients, so those hypotheses do not
literally include the present model. At exactly $s=3$, all-order
logarithmic corrections already appear in the marginal reports and the
integrated Gamma-moment chapter \cite{repo-marginal,repo-moments}.

We prove the missing joint theorem with every fixed logarithmic order,
including noncentral locations. Its Fourier remainder is uniform over
the whole lattice. A further argument converts that absolute control
into a relative result in a specified expanding positive window.
The latter is a partial positive-side resolution of Q2 in
\cite{repo-boundary}; it does not cover arbitrary deviations or the
negative saddle window.

\subsection{What comes from classical analysis}

The convergent Jonqui\`ere expansion, Gamma recurrence, and
polylogarithm continuation are classical \cite{dlmf-polylog}.
Lagrange inversion and its probability interpretation supply the exact
coefficient representation already used in the repository. Neither
these tools nor Gaussian local limits are claimed as discoveries.

There is also substantial prior work on nonclassical corrections to
Gaussian limits. Lam, Blanchet, Burch and Bazant \cite{lam}
express heavy-tail corrections using parabolic cylinder functions and
discuss logarithmic effects and matching to the tails. Berger, Birkner
and Yuan \cite{berger} study Gaussian and single-large-jump local
deviations, including the critical tail index two. Thus the appearance
of parabolic cylinder functions and the Gaussian--power-tail transition
are established phenomena. Their occurrence here is useful because
the polylogarithmic order moves with $n$ through the collision of two
singular terms, and the nonlinear feedback is retained exactly.

The proposed contribution is the combined formulation: a uniformly
pole-paired scale, a finite all-order Stirling calculus for a moving
order and moving feedback, a fully proved uniform remainder, and the
two-derivative composition principle. The local ramification theorem
later in the article is another direct deduction from the classical
polylogarithm expansion, with a precise repository application. The
limited comparison performed here cannot certify first priority in
the entire literature.

\subsection{Results at a glance}

\begin{center}
\begin{tabular}{@{}>{\raggedright\arraybackslash}p{0.22\textwidth}>{\raggedright\arraybackslash}p{0.69\textwidth}@{}}
\toprule
Result & Mathematical content\\
\midrule
Exact normalization & A unique scale $B>1$ solves a paired
Gamma--zeta equation; it is a Lambert expression when $s=3$.\\
All fixed orders & A finite sum of Gaussian order derivatives gives
error $O(L^{-N-1})$, uniform in lattice location.\\
Central identities & Every central correction is a finite Gamma
moment; the marginal coefficients match the previous ones exactly.\\
Composition principle & Higher endpoint derivatives of an analytic
outer function enter only beyond all fixed inverse-logarithmic orders.\\
Expanding positive window & Relative power-tail matching holds up to
any fixed power of $L$; the transition has an explicit Lambert center.\\
Inverse endpoint theorem & Minimal finite ramification is classified
for every real $s>1$, including all removable reciprocal orders.\\
\bottomrule
\end{tabular}
\end{center}

\subsection{Notation and the meaning of the asymptotics}

All moving-order theorems concern real $\eps$ with
$|\lambda|\le\Lambda$ for an arbitrary fixed $\Lambda<\infty$.
The positive feedback coordinate $b$ remains in a fixed compact interval
$[b_-,b_+]\subset(0,\infty)$. Constants in an $O$ term may depend on
these sets and the fixed truncation order, but not on $n$, the parameter
choice within them, or the lattice location when uniformity in that
location is asserted. A statement to every fixed order does not permit
the truncation order to grow with $n$.

The logarithm $\Log t$ is principal on the slit plane. On the Fourier
axis,
\begin{equation}\label{eq:fourier-log}
 \Log(-iy)=\log|y|-\frac{\pi i}{2}\operatorname{sgn}(y)
 \qquad(y\in\R\setminus\{0\}).
\end{equation}
The normal density is
$\normal(\xi)=(2\pi)^{-1/2}e^{-\xi^2/2}$.
The notation $\HH_a^{[j]}$ always means differentiation in the
\emph{order} $a$, not in the spatial variable $\xi$.
The coefficients $\stir{j}{k}$ are Stirling numbers of the second kind.

Inverse-logarithmic sectors in $L$ are transseries sectors in the usual
asymptotic sense. The nonlinear local remainder below is of size
$n^{-c}$ times logarithms, hence of exponential size $e^{-cL}$ relative
to those sectors. We identify its scale separation but do not claim
to compute all such exponential sectors.
